\section{数据集构建}
\label{sec:dataset-construction}

图~\ref{fig:pipeline}展示了如何将每个多轮样例从一条成功的强模型轨迹逐步降级，并通过经执行验证的搜索生成每个步骤的已验证层级标签。

\begin{figure}[t]
\centering
\includegraphics[width=\linewidth]{figures/fig2.pdf}
\caption{TwinRouterBench 构建流水线。对于每个多轮样例，流水线从成功的强模型轨迹出发，经由执行验证的搜索将各个步骤逐步降至更低成本的层级，从而为轨迹中的每次 LLM 调用生成已验证的层级标签。}
\label{fig:pipeline}
\end{figure}

\paragraph{从成功的强模型轨迹开始。}
我们首先收集在强模型配置下得到解决的轨迹。对于 SWE-bench，通过 \texttt{mini-swe-agent} 运行的 Claude Opus 4.6 或 GPT-5.4 必须提交能通过官方 \texttt{FAIL\_TO\_PASS} 测试的补丁；BFCL、mtRAG、QMSum 和 PinchBench~\citep{pinchbench} 使用各自 harness 的通过谓词。丢弃失败的种子轨迹，以避免路由标签将任务难度与路由错误混为一谈。

\paragraph{在因果前缀下搜索更低层级。}
对于每条保留下来的轨迹，Claude Opus 4.6 提供一个搜索提示：该步骤看起来是否可降级，以及若可降级，应首先探测的最激进低层级 $h_i$。该提示只是剪枝工具，绝非标签。随后我们运行贪心的顺序锁定降级搜索（附录算法~\ref{alg:seq-lock}）。在第 $i$ 步，先前锁定的步骤保持其已验证层级，未来步骤使用强层级，并从 $h_i$ 开始，按由低到高的顺序尝试第 $i$ 步的候选。因为逐步正确性并非总能直接观测，只有当完整轨迹仍以相同步数通过时，才接受候选，并推断降级步骤被正确处理；若无候选通过，该步骤保留在强层级。该方法将朴素的 $|T|^N$ 网格降为 $\mathcal{O}(|T|N)$ 次试验，同时确保后续前缀包含先前路由步骤实际产生的降级输出。锁定前缀由以实例、步骤、模型和生成参数为键的响应缓存重放；当上游锁定项改变时，缓存条目失效。

\paragraph{以模型池而非单一模型验证层级。}
层级标签应表示该层级的模型池能够解决该步骤，而不应表示某一个模型恰好通过。对于非 \texttt{low} 标签，我们运行三模型级联：只有该层级池中任一模型在混合模型轨迹下通过，才接受该层级。在 SWE-bench 最后一步审计中，级联在 33 个可降级标签中恢复 28 个，而单模型探测仅恢复 15 个，这证实单一代表模型对监督而言过于脆弱。

\paragraph{加固开放式通过谓词。}
一些基准（mtRAG、QMSum）依赖 LLM-as-judge 评测，其默认提示词和评分准则不够严格。在早期开发中，我们人工检查了随机样本的裁判判决，发现表面相似掩盖了错误或不完整答案的伪通过。我们以结构化裁判替换这些弱的参考相似度准则：先检查当前轮是否已解决，再评估忠实性、恰当性和完整性；并对证据冲突设置硬失败、在不确定时保守地失败。我们根据 Radbench 风格的阈值评分校准加固后的准则~\citep{kuo2025radbench}，并通过迭代测试验证其消除了初始样本中发现的伪通过模式；结构上不可靠的 mtRAG 样例和所有层级均失败的样例被丢弃。SWE-bench 与 BFCL 已具有可执行或结构化的通过谓词，无需额外加固。

\paragraph{人工审计。}
我们的降级搜索固定其余全部步骤，并一次测试一个步骤。该贪心策略将计算成本从指数级降至线性级，但可能遗漏跨步骤的交互效应：两个降级的组合可能失败，即使它们各自单独成功。为防范此类情形，我们对具有多步轨迹的三个工作负载（BFCL、SWE-bench、PinchBench）开展人工抽样审计，从每个工作负载路由步骤中独立随机抽取约 ${\sim}10\%$。mtRAG 和 QMSum 是缺少多轮累积复杂度的单轮任务；其标签已受加固裁判约束，加固过程已通过迭代校准消除了伪通过模式（\S\ref{sec:dataset-construction}），因此未在最终审计中重新抽样。对于每个抽样步骤，审阅者固定路由器可见的前缀，并判断该层级是否还能再降一级，将其标记为\emph{紧致}、\emph{不确定}或\emph{可继续降级}。审计覆盖 64 个步骤（25 个 BFCL、34 个 SWE-bench、5 个 PinchBench）：其中 63 个被判定为紧致；一个 SWE-bench 步骤被标记为可继续降级，且其层级在发布前已下调（表~\ref{tab:manual-audit-main}）。审计协议细节见附录~\ref{app:manual-audit}。

\begin{table}[t]
\centering
\small
\begin{adjustbox}{max width=\linewidth}
\begin{tabular}{lrrrrr}
\toprule
基准 & 总步骤数 & 抽样步骤数 & 抽样比例 \% & 可降级 & 紧致比例 \% \\
\midrule
BFCL & 248 & 25 & 10.1 & 0 & 100.0 \\
SWE-bench & 336 & 34 & 10.1 & 1 & 97.1 \\
PinchBench & 48 & 5 & 10.4 & 0 & 100.0 \\
\midrule
总计 & 632 & 64 & 10.1 & 1 & 98.4 \\
\bottomrule
\end{tabular}
\end{adjustbox}
\caption{步级人工审计汇总。\emph{可降级}表示审阅者认为已验证层级还可再降低一级而不破坏任务解决；\emph{紧致比例}表示已处于条件最优层级的比例。我们从具有多步轨迹的三个工作负载中抽取路由步骤；mtRAG 和 QMSum 是已由加固裁判保护的单轮任务，未重新抽样。}
\label{tab:manual-audit-main}
\end{table}

\paragraph{发布准则。}
只有同时满足以下条件的记录才能进入静态赛道：(1) 强模型种子轨迹通过；(2) 层级池中存在一个模型在混合轨迹下通过；(3) 对于开放式工作负载（mtRAG、QMSum），加固裁判不将答案标记为伪通过。每个标签都是能够正确处理当前步骤的最低成本层级的执行验证估计：由于直接的逐步判断并非总是可能，我们从固定步数下的端到端轨迹通过推断逐步充分性，并通过人工审计交叉检查。
